\subsection{Definición estructural y función de la constante gravitatoria}
\label{sec:ii-definicion-gravedad}

\begin{definition}[Sección gravitatoria contragrediente de la acción]
Para una sección positiva \(\hbar_a\), una velocidad
\(c_{\rm int}>0\) y el radio cronológico \(R_{108}\) de la
sección~\ref{sec:gravity}, la sección gravitatoria circular es
\begin{equation}
 G_a=\frac{c_{\rm int}^{3}R_{108}^{2}}{\hbar_a},
 \qquad R_{108}=\frac{54}{\pi}\ell_0,
 \qquad \ell_0=c_{\rm int}t_0.
 \label{eq:ii-definicion-gravedad-seccion}
\end{equation}
Es la sección recíproca de la acción que realiza la coincidencia
entre el radio gravitatorio reducido del ciclo y su radio
cronológico, en la convención \(r_g=Gm/c_{\rm int}^2\).
\end{definition}

La proposición de coincidencia de radios de la
sección~\ref{sec:gravity} demuestra su unicidad en la colección
\(G_a(\vartheta)=\vartheta c_{\rm int}^{5}t_0^2/\hbar_a\).
La condición circular fija \(\vartheta=(54/\pi)^2\) y conduce
a \eqref{eq:ii-definicion-gravedad-seccion}. Esa condición y las
secciones dimensionales positivas forman parte de la realización.
La sección~\ref{md:rigidez} construye su longitud y reúne la
prueba de unicidad para el carácter positivo general, con
$R_{108}=L_\alpha$. Por ello el selector circular y el lector
longitudinal conservan una misma normalización.
Manteniéndolas comunes entre las dos orientaciones, se obtiene
\begin{equation}
 \frac{G_{\mathrm{pre}}}{G_{\mathrm{ret}}}
 =R_{\mathrm{act}}^{-1},
 \qquad
 \frac{G_a\hbar_a}{c_{\rm int}^{3}}
 =R_{108}^{2}=\ell_P^2.
 \label{eq:ii-definicion-gravedad-area}
\end{equation}
La prueba es la sustitución directa de
\eqref{eq:ii-definicion-gravedad-seccion}; el cambio de acción
se compensa por el cambio recíproco de \(G_a\).
